% Source presentations for register row N-c-007.

\subsection{Two guarded-fragment presentations}
\label{sec:p09-presentation-n-c-007}

\subsubsection{Andr\'eka--van Benthem--N\'emeti's modal genealogy}

The registered article is the 1998 \emph{Journal of Philosophical Logic}
publication, but the technical manifestation verified for this reconstruction
is the authors' 1996 ILLC preprint.  All section and theorem locators in this
subsection refer to that preprint; they are not asserted to be the pagination
or final wording of the journal version \cite{AndrekaVanBenthemNemeti1998}.

The preprint compares three bounded quantifier forms.  In schematic binary
notation they permit
\[
 \exists y\,(Ryx\land\varphi(y)),
 \qquad
 \exists y\,(Ryx\land\varphi(x,y)),
 \qquad
 \exists y\,(Ryx\land\varphi(x,y,z)),
\]
with the corresponding tuple and higher-arity versions.  Formulas start from
atoms and close under Boolean connectives and the permitted quantifiers.  In
the middle fragment, called the \emph{guarded fragment}, one atomic guard must
contain every variable free in the quantified matrix.  Equality atoms may
occur in formulas but are not themselves guards, and Boolean combinations of
atoms are not guards \cite[\S4.1]{AndrekaVanBenthemNemeti1998}.  The first
fragment stays closest to the modal standard translation, the guarded fragment
can express symmetry but not transitivity of a binary relation, and the third
fragment is already undecidable.  Theorem~4.2.3 makes the three an expressive
proper hierarchy.

A finite set of elements is \emph{guarded} when it is exactly the set of
entries in some true relation tuple.  A guarded bisimulation is a nonempty
family of finite partial isomorphisms with back-and-forth extensions to guarded
sets, each extension agreeing with the current map on the overlap.  These maps
inspect relational information on guarded tuples rather than arbitrary finite
substructures.  Proposition~4.2.1 proves invariance of guarded formulas;
Theorem~4.2.2 proves the converse, that a first-order formula invariant under
guarded bisimulation is equivalent to a guarded formula.  Its finite-round
version characterizes the corresponding bounded-variable fragment
\cite[\S4.2]{AndrekaVanBenthemNemeti1998}.

Theorem~4.6.1 proves decidability for the guarded fragment \emph{without
identity}.  For the input formula \(\varphi\), the proof forms a finite set
\(F\) of bounded-length formulas in its variables, closed under the required
Boolean and substitution operations.  An \(F\)-type is a maximally consistent
choice from that finite set.  A \emph{quasi-model} is a finite set of such
types satisfying the source's successor condition for every guarded
existential formula.  Every ordinary model yields a quasi-model; conversely,
paths through a quasi-model are unfolded into an ordinary model, and a truth
lemma transfers the formulas in \(F\).  Existence of the finite quasi-model is
effectively checkable, which yields the decision procedure
\cite[\S4.6]{AndrekaVanBenthemNemeti1998}.

The word ``finite'' in that proof modifies the abstract quasi-model.  The
unfolded ordinary model may be infinite, and the inspected preprint explicitly
says that the argument does not establish completeness with respect to finite
ordinary models.  It assigns the guarded finite-model property to later work
and warns that equality requires additional syntactic care.  Because the final
1998 article was not available for clause-by-clause comparison, this preprint
boundary is not used to assert what revision occurred in the journal version.

\subsubsection{Gr\"adel's model and complexity boundaries}

Gr\"adel works over relational vocabularies, allowing equality and constants
but no positive-arity functions.  The guarded fragment closes atoms under
Boolean operations and quantifiers of the form
\[
 \exists\bar y\,
 \bigl(\alpha(\bar x,\bar y)\land
       \varphi(\bar x,\bar y)\bigr),
\]
where the single atomic guard \(\alpha\) contains every free variable of the
matrix.  The \emph{loosely guarded fragment} permits a conjunction of atoms as
guard, provided each newly quantified variable co-occurs in some conjunct
with every variable that remains relevant.  The second condition is strictly
weaker than one all-variable atom and is used to express temporal patterns
which plain guarded syntax cannot express \cite[defs.~2.1--2.2]{Gradel1999}.

Theorem~3.7 constructs, for every satisfiable loosely guarded sentence using
\(m\) variables, a bounded-degree tree-like model by normalizing the sentence
and repeatedly attaching witness structures while preserving their atomic
types.  This model may be infinite.  For the ordinary guarded fragment,
Theorem~3.10 proves the finite-model property: choose a finite representative
subtree for the finitely many local isomorphism types, then apply Herwig's
extension theorem to close the partial isomorphisms while retaining every
guarded witness obligation \cite[\S3]{Gradel1999}.  The paper leaves the
finite-model property of the loosely guarded fragment open; the two fragments
must not share one unqualified finite-model row.

The complexity proof also has two sides.  Theorem~4.3 decides loosely guarded
satisfiability in deterministic double-exponential time by searching through
exponentially represented atomic types.  Theorem~4.4 encodes an alternating
exponential-space computation tree in guarded formulas, proving
\(2\mathsf{EXPTIME}\)-hardness already for the guarded fragment.  Hence guarded
and loosely guarded satisfiability, with or without equality, are
\(2\mathsf{EXPTIME}\)-complete.  If either the number of variables or the
maximum relation arity is fixed, the classification falls to
\(\mathsf{EXPTIME}\)-complete; over a fixed finite vocabulary the loosely
guarded problem is in \(\mathsf{EXPTIME}\)
\cite[thms.~4.3--4.4 and cors.~4.5--4.6]{Gradel1999}.

Finally, the paper proves that adding functionality statements or stipulated
transitive relations crosses an undecidability boundary by finite-grid
encodings \cite[thms.~5.8 and 5.10]{Gradel1999}.  These statements concern
extensions of the language.  They do not mean that functionality or
transitivity was definable in the base guarded fragment, and they do not alter
the modal-standard-translation result in
Theorem~\ref{thm:p09-modal-guarded}.
